\section{Introduction}
For a vertex $v$ of a finite oriented graph $D$, let $N_D^+(v)$ denote its
outneighbourhood and let $N_D^{++}(v)$ consist of vertices whose directed
distance from $v$ is exactly two. Seymour's conjecture asks whether every
nonempty oriented graph has a vertex satisfying
\begin{equation}\label{eq:paper-snc}
 |N_D^{++}(v)|\ge |N_D^+(v)|.
\end{equation}
Such a vertex is called a Seymour vertex. The exclusion of first neighbours
from the second set is essential in every deletion, completion and counting
argument below.

The tournament case is known, but a general completion may create second
neighbours that were absent in the original graph. We develop two ways of
making that discrepancy explicit. The first uses a small induced incoming
core and an exact linear certificate. The second completes only carefully
controlled directions before a global median-order optimization. A separate
copying operation imposes global structural restrictions on a suitably chosen
counterexample. The three constructions have different preservation properties;
their hypotheses are kept distinct throughout the paper.

\subsection{Principal results}
Partition vertices by equality of their full incoming neighbourhoods. A class
is maximal when its common incoming set is maximal under inclusion among the
sets $N_D^-(v)$. Maximality refers to inclusion, not to the numerical indegree.
Our first unconditional result is the following computer-assisted theorem.

\begin{theorem}[Five maximal incoming classes]\label{thm:intro-five}
Every finite nonempty oriented graph with at most five maximal incoming-neighbourhood
classes has a maximal class all of whose vertices are Seymour vertices.
The total order and the sizes of the classes are unrestricted.
\end{theorem}

The theorem follows from a transfer inequality valid for arbitrary induced
cores. Its finite component comprises $59\,809$ labelled cores of orders one
through five and $462\,404$ nonempty pattern inequalities. Every witness is a
nonnegative integer vector, with entries at most three. The tables are checked
independently of the program that found them. The full proof and the precise
finite verification contract appear in Theorem~\ref{thm:core-five}.

The same certificate system has a converse: an infeasible core system yields,
by a finite labelled construction, a genuine all-deficient oriented graph
(Theorem~\ref{thm:pattern-duality}). Thus universal feasibility is equivalent
to the full conjecture. A certificate also lifts through repeated maximal
incoming-core reductions, reducing the remaining feasibility question to
cores with pairwise distinct, incomparable incoming sets.

Our second group of results concerns simultaneous copying. If an original
vertex map $\rho$ satisfies $N_D^-(v)\subseteq N_D^-(\rho(v))$ for every $v$,
replacing each outgoing row by the original row of $\rho(v)$ creates an
oriented graph whose gap at $v$ is at least the original gap at $\rho(v)$.
This implies a global covering property in a smallest-order counterexample:
every vertex receives an entire maximal incoming class whose members all have
gap one. A further extremal selection makes equal incoming neighbourhoods into
independent true-twin modules. These statements use global extremality rather
than merely inclusion-minimal deletion assumptions.

For the completion results write $g_D(v)=d_D^+(v)-d_D^{++}(v)$ and let $t_D(v)$
count the original far vertices not protected by incoming-neighbourhood nesting.
The closure constructed in Section~\ref{sec:paper-closure} admits an optimal
reference completion whose every original-positive-gap median feed satisfies
\begin{equation}\label{eq:intro-surplus}
 t_D(f)\ge g_D(f)+1.
\end{equation}
Under the safe weak extension, zero completed gap requires at least three
remaining missing partners, an original directed cycle among them, and
$t_D(f)\ge g_D(f)+3$. For a strongly connected all-deficient graph minimal
under every nonempty arc-set deletion, the same feed satisfies
$\eta_D(f)\ge4$, where
$\eta_D(v)=\mu_D(v)-1+d_D^{--}(v)-d_D^{++}(v)$.

The supplementary case reductions, together with an unbounded one-exception
theorem for actual bad-direction failure sources, yield
\begin{equation}\label{eq:intro-residual}
 2M-n\ge7,\qquad
 M\ge\left\lceil n/2\right\rceil+3+\ind_{\{n\text{ even}\}}.
\end{equation}
The missing-pair bound is a necessary condition for the stated arc-set-minimal
class. It is not a missing-edge density theorem for arbitrary original graphs:
deleting arcs changes $M$. Proofs of the residual allocations and the specialized
tools for $n=2\delta^++3$ are collected in the companion supplement.

\subsection{Relation to previous work}
\label{subsec:related-work}

The tournament case of Seymour's second neighborhood conjecture was proved
by Fisher~\cite{Fisher1996} using a probability distribution obtained from
a skew-symmetric game. Havet and Thomass\'e~\cite{HavetThomasse2000}
subsequently gave a proof using median orders, including a stronger bound
involving the good vertices of an order. Both methods enter the present
work. The incoming-core argument uses the existence of a losing density
on a tournament; the completion arguments use the strong feed bound and
the equality case of sedimentation. These are established inputs. Our
arguments specify the additional conditions needed to transfer information
from a tournament completion to the exact second neighborhoods of the
original graph. In particular, the existence of a Seymour vertex in a
completion alone does not provide such a transfer.

Fidler and Yuster~\cite{FidlerYuster2007} proved the vertex-weighted
tournament statement and established further cases obtained by deleting
prescribed edge sets from tournaments. Ghazal's work on generalized stars
and other missing graphs~\cite{Ghazal2012,Ghazal2013} develops the convenient
orientation and good-missing-edge approach underlying several of our
completion definitions. Dara, Francis, Jacob and
Narayanan~\cite{DaraEtAl2022} extended the matching case to a missing graph
whose edges can be partitioned into a matching and a star, and also
considered partitions into an independent set and a $2$-degenerate graph.
Our protection-preserving completions require additional compatibility
conditions beyond convenience of individual added arcs. The proof of each
transfer inequality therefore accounts for paths using added arcs and for
original second neighbors that become first neighbors. None of these
inequalities is inferred simply by applying a known tournament theorem to
an arbitrary completion.

Minimal-counterexample methods provide another starting point.
Brantner, Brockman, Kay and Snively~\cite{BrantnerEtAl2009} established
strong connectivity, gaps in $\{1,2\}$, and restrictions involving
gap-one predecessors and directed cycles for their minimal counterexamples.
Their extremal selection first minimizes the number of arcs over all
counterexamples and then minimizes the order. The selections used here
must be distinguished from that convention: fixed-order minimality under
deletion of a nonempty arc set, and minimum order followed by arc count
and further finite objectives, serve different arguments. We give the
required deletion and copying arguments for the stated selection in each
case. In particular, the predecessor-containment row operation and the
resulting normal form are not treated as consequences of the earlier
extremal definition.

Seacrest's subset formulation~\cite{SeacrestSubsets2019} supplies a closely
related boundary perspective. Its edge-minimal subset inequality motivates
the boundary budgets used in the residual analysis. Because the sets and
minimality assumptions have to agree exactly, the external-boundary form
needed here is proved explicitly. This paper is distinct from Seacrest's
arc-weighted extension~\cite{Seacrest2015}; an arc-weighted result is not
substituted for the unweighted feed or boundary statements. The residual
bound $R\geq 7$ in this manuscript concerns the specified minimal
all-deficient graphs. It must not be read as an assertion about the
missing-edge count of every counterexample before a minimizing deletion.

There are further reasons to keep extremal assumptions visible.
Espuny D\'iaz, Gir\~ao, Granet and
Kronenberg~\cite{EspunyDiazEtAl2025} prove that a vertex-minimal
counterexample has minimum outdegree greater than the square root of its
order, while also constructing arbitrarily large strongly connected
counterexamples with bounded minimum outdegree, conditional on the
existence of any counterexample. Zelenskiy, Darmosiuk and
Nalivayko~\cite{ZelenskiyEtAl2021} obtain conditional counterexample
constructions with widely differing densities and diameters. Thus
counterexample amplification, and cyclic substitution in particular,
already have a literature. The calculations here involving such
constructions are used to delimit particular completion or weighted
transfer claims, rather than to assert a new general amplification
principle.

Theorem~\ref{thm:core-five} has a different parameter from these density
and degree results: the number of distinct inclusion-maximal incoming
neighborhoods. There is no restriction on the order of the original graph
or on the size of an incoming-neighborhood class. The certificate on a
chosen core is checked against every permitted subset of an incoming set,
so that it controls all original target patterns simultaneously. An
ordinary verification of Seymour's conjecture for small graphs does not
automatically supply these stronger certificates. The finite tables are
therefore accompanied by explicit integer checks and a proof of their
transfer to arbitrary-order graphs. The labelled construction from an
infeasible pattern system establishes a converse at the universal level.
Finite linear-programming duality is standard; the graph realization and
its exact-distance calculation are the substantive steps that must be
verified for this particular reformulation. Universal feasibility remains
unproved.

Recent special cases provide useful comparisons. Wang and
Lu~\cite{WangLu2026} prove the conjecture for
$\mathcal D_{0,3}\cup\mathcal D_{1,1}$, where $\mathcal D_{s,t}$ consists
of oriented graphs admitting a vertex partition whose induced underlying
graphs are respectively $s$- and $t$-degenerate. Bai, Li and
Park~\cite{BaiLiPark2026} study the stronger requirement of a complete
directed matching from the first to the exact second neighborhood, proving
it when the minimum outdegree is at most five or the graph is
$5$-anti-transitive. Kaneko and Locke~\cite{KanekoLocke2001} established
ordinary Seymour vertices when the minimum outdegree is at most six.
The preprint of Sadhukhan, Sandeep and Sen~\cite{SadhukhanEtAl2026}
treats minimum outdegree seven using computer-assisted local reductions.
Brukhman's preprint~\cite{Brukhman2026} treats the dense range
$n\leq 2\delta^++2$. The incoming-core theorem uses none of these
recent preprints as a premise.

The five-class hypothesis does not imply the hypotheses of those
particular recent results. For example, replace each vertex of the regular
five-vertex tournament, with arcs $i\to i+1,i+2$ modulo five, by an
independent set of size $m\geq4$. The resulting graph has five maximal
incoming classes, order $5m$, and minimum outdegree $2m$. Its underlying
complete five-partite graph belongs to neither
$\mathcal D_{0,3}$ nor $\mathcal D_{1,1}$, and a path through parts
$0,1,2,3,4,1$ excludes $5$-anti-transitivity. It also lies outside the
stated dense range. This example itself is already covered by the weighted
tournament theorem. It demonstrates noncontainment in specific hypotheses,
not novelty relative to every known weighted or substitution argument.

Halkiewicz's split-twin operation~\cite{Halkiewicz2026} duplicates incoming
neighbors and partitions outgoing neighbors while increasing the order;
its stated preservation theorem has a condition involving an existing
Seymour vertex. This differs from the fixed-order row copying used here.
Nevertheless, this distinction and the preceding comparisons do not settle
priority for equivalent formulations under other terminology. The present
paper makes no claim to an exhaustive classification of earlier reductions.
The literature comparison is current to 7 October 2026. Public full-proof
claims, including~\cite{GloverClaim2026}, are not assumed correct or used
as premises. The conclusions proved here remain partial results and
reformulations of the conjecture.


\subsection{Organization and verification}
The incoming-core and cloning arguments are independent of the residual
machinery. After establishing them, we give the closure construction and the
zero-gap cycle argument, then state the minimal-counterexample consequences
with their supplementary proof locations. The final sections give reproducibility
details, exact controls against stronger transfer claims, and the unresolved
global steps. The supplement is a self-contained technical repository of the
longer foundational and case arguments. It also records which earlier finite
classifications are subsumed by the stronger results.
